\enteteExam{NFE106 -- Administration de bases de données}%
{Examen 11 septembre 2025 -- session 2 -- 3h}%
{Examen en présentiel -- Calculatrices autorisées -- Documents autorisés: 5 pages recto/verso de notes personnelles}%

Pour l'ensemble de cet examen, nous utiliserons le schéma relationnel suivant.

\begin{itemize}
 \item \textbf{Marque} (\underline{Id\_Marque}, Nom, siègeSocial)\\
 \item \textbf{Automobile} (\underline{Modèle, série}, id\_marque, annee\_sortie, prix\_base)\\
 \item \textbf{Option} (\underline{Modèle, série, id\_equipement})\\
 \item \textbf{Equipement} (\underline{Id\_Equipement}, Nom, Prix, Description)\\
\end{itemize}

Une voiture est désignée par son modèle (C5, 308, Megane\ldots) et sa série (Break, Sport, BZH\ldots).
Un modèle a 10 séries différentes en moyenne. Voici la taille des différents attributs :
\begin{itemize}
    \item Les identifiants (Id\_...) et les nombres (année, prix) sont codés sur 4 octets,
    \item Les attributs nom, modèle et série sont codés sur 50 octets
    \item Les attributs siègeSocial et description sont codés sur 200 octets,
\end{itemize}

Statistiques :
\begin{itemize}
    \item 135 nuplets dans la table {\tt Marque},
    \item 120 000 nuplets {\tt Automobile},
    \item 2 600 000 nuplets {\tt Option},
    \item 8 100 nuplets {\tt Equipement}
\end{itemize}
%Ne pas oublier de vérifier le nombe de tables !!!!

Une adresse physique occupe 10 octets, nous disposons de 21 blocs en 
mémoire centrale, et un bloc fait 8 KO.

\section{Stockage et Indexes (8 points)}


\begin{enumerate}
    \item Présenter sous forme de tableau la taille (en nombre de blocs) prise par chaque table~{\tt (2~points)}
\begin{solution}
    La taille d'une bloc utile est de 7372 octets.\\
    \begin{tabular}{|l|l|l|l|l|}
    \hline
    \textbf{Table} & {\bf Taille nuplets}    & {\bf Nb nuplets par bloc} & {\bf Nb blocs}\\
    \hline
    \hline
    Marque      & 260o  & 27    & 5\\
    Automobile  & 120o  & 60    & 2 000\\
    Option      & 110o  & 65    & 40 000\\
    Equipement  & 265o  & 27    & 300\\
    \hline
    \end{tabular}
\end{solution}
    \item La table {\tt Option} est de loin la table la plus lourde. Identifier d'où provient ce coût et proposer 
       une évolution du schéma  qui réduirait qui réduirait considérablement le stockage 
  Donner également les impacts sur la table {\tt Automobile}. {\tt (2~points)}
\begin{solution}
    La taille provient de la clé primaire qui est composée de 2 chaines de caractère. Ainsi, 
	 en rajoutant un id dans la table {\tt Automobile}, en l'utilisant dans 
	  la table {\tt Option} nous n'avons plus que 2 entiers. Nous obtenons 2735 blocs.\\
    Concernant la table {\tt Automobile}, nous perdons 64 blocs pour stocker cet identifiant.
\end{solution}

    \item Réécrivez la requête suivante pour qu'elle soit adaptée au nouveau schéma :

\begin{minted}{sql}
    SELECT Serie, COUNT(*) 
    FROM Option 
    WHERE Modele = '308' 
    GROUP BY Serie;
\end{minted}
    Quel sera l'impact sur l'évaluation ? (une explication est demandée, pas un coût) {\tt (2~points)}
\begin{solution}
    \begin{verbatim}
    SELECT A.série, COUNT(*) FROM Automobile A, Option O
    WHERE A.ID=O.ID AND A.modele = '308' GROUP BY A.Serie;
    \end{verbatim}
    Le modèle étant reporté à un champ de la table {\tt Automobile} 
	 et n'étant plus présent dans la table {\tt Option}, 
    on ne peut plus faire de requête directement sur la table {\tt Option}. Du coup, 
    c'est une requête qui fera une jointure entre {\tt Automobile} et {\tt Option} 
    pour récupérer tous les id de voitures '308'. Cette jointure peut être un surcoût conséquent.
\end{solution}

  \item Quel est le nombre maximal d'entrées d'arbre B que l'on peut mettre dans un bloc,
     avant modification et après modification~{\tt (2~points)}
%    \item En déduire la hauteur minimale  d'un arbre B sur la clé primaire de la table {\tt Option} : 
%      avant modification et après modification~{\tt (2~points)}
\begin{solution}
\begin{itemize}
  \item Avant : clé de 104o + 10o, 71 cases, ordre 35 $\Rightarrow$ hauteur 5
  \item Après : clé de 8o + 10o, 454 cases, ordre 227 $\Rightarrow$ hauteur 3
\end{itemize}
\end{solution}
\end{enumerate}

\section{Statistiques et Optimisation (8 points)}

\begin{enumerate}
  \item ({\tt 2~points})
  Pour la requête ci-dessus, donner la cardinalité de l'attribut {\tt id\_marque}
  et  donner le nombre de blocs à lire en supposant un index sur {\tt id\_marque} 


\begin{minted}{sql}
SELECT COUNT(*) FROM Automobile WHERE id_marque = 5
\end{minted}

  \item ({\tt 2~points}) Même question avec la requête 
  
\begin{minted}{sql}
SELECT modele, serie FROM Automobile WHERE id_marque = 5
\end{minted}

\begin{solution}
\begin{enumerate}
  \item Cardinalité : 1/130.
  \item Arbre B :  ordre 291, hauteur 3. Nous récupérons 1 000 ROWIDs, mais sans faire d'accès. Il faut donc compter les ROWIDs, en parcourant les feuilles de l'index. 4 feuilles à parcourir $\Rightarrow$ 3 + (4 - 1) = 6 I/O
\end{enumerate}
\end{solution}

  \item ({\tt 2~points}) Donner, pour le plan d'exécution ci-dessous, la requête SQL correspondante. 
   Donner également le sens de cette requête. Le modèle de la voiture sera projeté. 
\begin{verbatim}
0   SELECT STATEMENT
1     NESTED LOOPS
2       NESTED LOOPS
3         NESTED LOOPS
4*          TABLE ACCESS FULL            Marque
5*          TABLE ACCESS BY INDEX ROWIDS Automobile
6*            INDEX RANGE SCAN           IDX_Id_Marque_Automobile        
7*        INDEX RANGE SCAN               PK_Option
8*      TABLE ACCESS BY INDEX ROWIDS     Equipement
9*        INDEX UNIQUE SCAN              PK_Equipement

(4) Nom = 'Peugeot'
(5) Série = 'BZH'
(6) M.Id_Marque = A.Id_Marque
(7) A.Modele = O.Modele AND A.Serie = O.Serie 
(8) Nom = 'Wifi'
(9) O.Id_equipement = E.id_equipement
\end{verbatim}
\begin{solution}
\begin{verbatim}
SELECT A.NOM
FROM Marque M, Automobile A, Option O, Equipement E
WHERE  M.id_marque = A.id_marque AND
       A.Modele = O.Modele AND A.Serie = O.Serie AND
       O.id_equipement = E.id_equipement AND
       M.Nom = 'Peugeot' AND
       A.Serie = 'BZH' AND
       E.Nom = 'Wifi';
\end{verbatim}
Donner le nom des modèles de Peugeot de série 'BZH' qui ont un Wifi de série. 
\end{solution}

    \item ({\tt 2~points}) Pour les opérateurs des lignes 4 à 6, donner le nombre de nuplets en entrée et en sortie, sachant que :
\begin{itemize}
  \item Les noms de marque et d'équipement sont uniques,
  \item Les voitures sont uniformément réparties sur les marques
\end{itemize}
\begin{solution}
\begin{itemize}
  \item 1 tuple Peugeot accède à l'index pour obtenir $\frac{130 000}{130} =  1 000$ automobiles,
  \item $1000 \times 1\% = 10$ modèles correspondent,
  \item $\frac{2800000}{130000} = 22$ options par voiture. Donc 220 options pour celles qui nous intéressent,
  \item On sélectionne un équipement parmi 480 000. Nous devrions obtenir 0, tout au plus 1 modèle.
\end{itemize}
\end{solution}
\end{enumerate}


\section{Concurrence (4 points)}

Soit l'exécution concurrente suivante :

$$H = r_2[y] w_1[x] r_1[x] r_3[z] w_2[z] w_3[x] c_1 w_3[z] c_3 r_2[x] c_2$$


\begin{enumerate}
 \item Donner la liste des conflits de $H$. 

\begin{solution}

Conflits:
\begin{enumerate}
 \item En $x$: $w_1w_3$, $w_1r_2$, $r_1w_3$, et $w_3r_2$
  \item En $y$: -
  \item En $z$: $r_3w_2$, $w_2w_3$
\end{enumerate}

\end{solution}

 \item Donner le graphe de sérialisabilité de $H$. Que pouvez-vous déduire de ce graphe ? 

\begin{solution}
Les arêtes\,: $T_1 \to T_3$ ; $T_1 \to T_2$ ; $T_3 \to T_2$ ; $T_3 \to T_2$  et  $T_2 \to T_3$.

Il y a un cycle donc l'exécution n'est pas sérialisable.
\end{solution}

\item  Donner l'exécution obtenue 
par application de l'algorithme de verrouillage à deux phases. 
Détailler le déroulement de l'algorithme. 


\begin{solution}
La lecture $w_2[z]$ est bloquée par $r_3[z]$. $T_2$ est mise en attente et $T_3$ termine. $T_2$
est alors débloquée et termine.

$$H_{finale} = r_2[y] w_1[x] r_1[x] r_3[z] c_1 w_3[x] w_3[z] c_3 w_2[z] r_2[x] c_2$$
\end{solution}
\end{enumerate}



